\documentclass{beamer}
%Language
\usepackage[dutch]{babel} % use dutch hyphenation etc.
%math stuff
\usepackage{amssymb,amsmath,amsthm,dsfont,mathrsfs}

\usetheme{Warsaw}


\usepackage{listings}

% Set slide numbers
\setbeamertemplate{footline}[frame number]

\usepackage{hyperref}
%Document information
\title{\TeX niCie presentation} 
\author[Jagna Hendriks]{Jagna Hendriks}
\date{\today}

\begin{document}
	\begin{frame}{Tekstopmaak}
		\begin{tabular}{c p{3cm} c c p{3cm}}
			Resultaat & Code & & Resultaat & Code \\
			\hline
			\textbf{Tekst} & Tekst & & {\tiny Tekst} & Tekst \\
			\textit{Tekst} & tekst & & {\small Tekst} & Tekst \\
			\textsc{Tekst} & tekst & & {\large Tekst} & Tekst \\
			\underline{Tekst} & tekst & & {\huge Tekst} & Tekst \\
			\hline
			$\mathbf{Tekst}$ & ... & & & \\
			$\mathbb{T}$ & ... & & & \\
			$\mathcal{T}$ & ... & & & \\
			$\mathfrak{T}$ & ... & & & \\
			$\boldsymbol{\alpha}$ & ... & & &\\
		\end{tabular}
	\end{frame}
	
	\begin{frame}{Packages}
		Packages staan in de \textbf{preamble} van een document om meer commando's in het document te kunnen gebruiken. 
		Handige packages zijn onderanderen:
		\begin{itemize}
			\item Voor wiskunde: amssymb, amsmath, amsthm, mathrsfs
			\item Voor taal: babel, inputenc, fontenc
			\item Voor afbeeldingen: graphicx, float
			\item Voor lijsten: enumitem
			\item Voor linkjes: hyperref, cleveref
			\item Voor layout: geometry, fancyhdr
			\item Voor zelf afbeeldingen maken: tikz
		\end{itemize}
	\end{frame}
	
	\begin{frame}{Zelf oplossingen vinden}
		\LaTeX biedt veel opties, wat zowel handig als frustrerend kan zijn. Daarom is het belangrijk om te leren zoeken naar wat je wil bereiken in je document. 
		
		Het helpt om duidelijke zoekopdrachten te hebben, zoals bijvoorbeeld 'tabellen maken latex', 'citing in latex', 'error ... latex'. 
		
		Handige websites om te bezoeken:
		\begin{itemize}
			\item \url{https://tex.stackexchange.com/}
			\item \url{https://www.overleaf.com/learn/latex/Learn_LaTeX_in_30_minutes}
			\item \url{https://nl.wikibooks.org/wiki/LaTeX/Wiskundige_formules}
		\end{itemize}
	\end{frame}
	

\end{document}